\section{ROM 布局、复位代码与入口不变量}

\topic{入口先建立确定状态}
入口依次执行 \code{cli}、\code{cld}，清零数据段、附加段和栈段，设置栈顶，
再初始化串口。关闭中断是因为 IDT 尚不存在；清除方向标志保证字符串读取向
高地址推进；显式设置栈是因为 \code{call}、\code{ret} 和
\code{push} 已经依赖它。

\topic{CS 覆盖用于读取高端 ROM}
固件代码和常量位于高地址 ROM。入口的 CS 隐藏基址仍指向 ROM 高区，而 DS
被清零并用于低地址 RAM。使用 \code{cs lodsb} 可以按当前代码段读取日志
字符串，避免假设 ROM 在低端还存在别名。

\topic{源码走读}
下面的清单直接引用配套工程，书与代码保持同源：

\lstinputlisting[
  language={[x86masm]Assembler},
  firstline=1,
  lastline=177,
  caption={16 位复位与串口入口}
]{../../../../source/firmware/src/reset_and_serial.asm}

\topic{ROM 审计}
宿主工具不通过“文件存在”判断成功。它检查镜像恰为 131072 字节，复位偏移的
操作码为 \code{E9}，解码位移后目标为 \code{0xF000}，入口以
\code{CLI/CLD} 开始。这是对最终字节的独立验证，不依赖汇编源码意图。

\topic{入口指令建立的逐项不变量}\par
\begin{osgridlongtable}{L{0.20\textwidth}L{0.28\textwidth}L{0.36\textwidth}}
\textbf{指令或动作} & \textbf{建立的状态} & \textbf{缺失时的后果}\\ \hline
\endfirsthead
\textbf{指令或动作} & \textbf{建立的状态} & \textbf{缺失时的后果}\\ \hline
\endhead
\code{cli} & \code{IF=0} & 外部中断可能在 IDT 建立前改变控制流\\ \hline
\code{cld} & \code{DF=0} & \code{lodsb} 可能向低地址移动并越过字符串起点\\ \hline
清零 \code{DS/ES/SS} & 低地址数据和栈采用已知段基址 & 复位残留状态进入地址计算\\ \hline
设置 \code{SP=0x7000} & 栈位于约定 RAM 范围 & \code{call} 把返回地址写入未知位置\\ \hline
初始化 UART & 建立 8N1 轮询发送协议 & 串口输出不能作为可靠证据\\ \hline
\end{osgridlongtable}

初始化顺序具有数据依赖。写 \code{SS} 与 \code{SP} 之前不能安全使用
\code{call}；设置 \code{DS} 之前不能按低端数据模型访问 RAM；清除方向标志
之前不能假设字符串指令递增。汇编列表的顺序不是排版选择，而是一条状态证明。

\topic{栈把控制流转换成内存状态}
\code{call} 在 16 位入口中先把返回 IP 压入栈，再跳到过程入口；
\code{ret} 从同一位置恢复 IP。栈向低地址增长，所以首次调用会把两个字节
写到 \code{0x6FFE} 和 \code{0x6FFF}。只设置 SP 而没有设置 SS，或者把
栈放入 ROM，都会使最普通的过程调用失效。

当前固件没有嵌套中断和局部栈对象，仍需为调用深度保留边界。后续进入 C++ 后，
参数传递、保存寄存器、对齐和局部对象会显著增加栈使用量，届时栈范围要成为
BootInfo 和内存保留表的一部分。

\topic{代码段覆盖的精确含义}
\code{lodsb} 默认从 \code{DS:SI} 读取一个字节到 AL，然后根据 DF 更新 SI。
前缀 \code{cs} 只把本次内存读取的段来源改为 CS，不改变 SI，也不永久修改
DS。复位近跳转保留了 CS 隐藏基址，因此标签相对 ROM 入口的偏移仍可用于
读取常量。

这一方案服务于最早入口，并不是通用数据模型。进入保护模式和长模式后，段的
含义变化，内核将依赖页表和链接地址访问只读数据。阶段性做法必须写清适用
条件，才能在条件消失时被安全替换。

\topic{源码、反汇编与执行轨迹}
源码展示程序员表达的符号和意图；反汇编展示最终地址上的机器指令；执行轨迹
展示处理器实际走过的控制流。三者应该一致，却不能互相替代。例如源码中一条
\code{jmp reset_entry} 会根据距离编码成短跳转、近跳转或带重定位的形式，
最终字节才是 CPU 接受的事实。

裸机调试从最终产物向上追溯：先确认 QEMU 取到的 ROM 字节，再确认链接地址，
最后回到源码。只盯着源文件会错过链接布局、填充和错误镜像参数。
